\chapter{The infinite fingerprint classification}
\label{ch:infinite-fingerprint}\label{chap:infinite-fingerprint-classification}

\begin{abstract}
The $G/L/C/M$ quadrichotomy of
Proposition~\ref{prop:three-invariants-relations}(iii) is a coarse
projection. It reads a single slot ($r_{\max}$) of a richer
invariant, and it degenerates on the critical-level locus
$\kappa_{\mathrm{ch}}(\cA) = 0$. The canonical fingerprint
$\varphi(\cA) = (p_{\max}, r_{\max}, \chi_{\mathrm{VOA}}, n_{\mathrm{strong}}, \mathrm{coset})$
of Definition~\ref{def:canonical-fingerprint} is the correct
complete invariant: it separates every pair of chiral algebras in the
standard landscape whose bar complexes differ, it survives
$\kappa_{\mathrm{ch}} = 0$ by recording the critical boundary stratum
$\mathrm{FF}$, and it transports functorially under
Drinfeld--Sokolov reduction. This chapter inscribes the four
strengthening theorems that complete the fingerprint programme: the
critical-level bar structure on class~$\mathrm{FF}$
(Theorem~\ref{thm:fifth-class-FF}); the pole--depth independence
table (Theorem~\ref{thm:pole-depth-independence}); the bijection
$d_{\mathrm{alg}}(\cA) = r_{\max}(\cA) - 2$
(Theorem~\ref{thm:d-alg-r-max-bijection}); and the DS transport law
(Theorem~\ref{thm:DS-fingerprint-transport}).
\end{abstract}

\section{Setup: the coarse projection degenerates at
\texorpdfstring{$\kappa_{\mathrm{ch}} = 0$}{kappa = 0}}
\label{sec:coarse-degeneration}

Chapter~\ref{ch:three-invariants} introduced three invariants
$(p_{\max}, k_{\max}, r_{\max})$ and the canonical fingerprint
$\varphi(\cA)$
\textup{(}Definition~\ref{def:canonical-fingerprint}\textup{)}. The
stratum theorem
\textup{(}Theorem~\ref{thm:five-class-stratum}\textup{)} partitioned
the standard landscape into the four generic shadow-depth classes
$G,L,C,M$ together with the critical boundary stratum~$\mathrm{FF}$
according to the pair~$(r_{\max}, \kappa_{\mathrm{ch}} = 0)$. Two points of that
theorem were stated but not proved in detail: first, the
internal structure of class~$\mathrm{FF}$, where the shadow metric
$Q_L$ degenerates and every shadow-based argument loses its pivot;
second, the exact relationship between $r_{\max}$ and the algebraic
depth~$d_{\mathrm{alg}}$. These gaps are the subject of this
chapter.

\begin{remark}[Why the quadrichotomy is insufficient]
\label{rem:why-quadrichotomy-insufficient}
The quadrichotomy $\{G,L,C,M\}$ is the image of the projection
$\Pi_{\mathrm{coarse}}\circ\varphi$ restricted to
$\kappa_{\mathrm{ch}}(\cA)\neq 0$. Three deficiencies follow. (i)
At~$\kappa_{\mathrm{ch}}=0$ the Sugawara construction
$T = (k+h^\vee)^{-1}:\!J^aJ^a\!:$ has a pole, the shadow metric
$Q_L(t) = (2\kappa_{\mathrm{ch}} + 3\alpha t)^2 + 2\Delta t^2$
loses its leading constant, and the shadow depth alone cannot
separate critical-level affine $V_{-h^\vee}(\fg)$ from
its~$\mathrm{FF}$ companions. (ii) Two algebras with identical
$r_{\max}$ but different parity (symplectic boson versus symplectic
fermion, $\Walg_3$ versus super-$\Walg_3$) share a quadrichotomy
class but not a bar complex. (iii) The DS reduction
$\cA\mapsto H^\bullet_{\mathrm{DS}}(\cA)$ scrambles the quadrichotomy
(affine~$\widehat{\fg}_k$ in class~$L$ maps to~$\Walg_N$ in
class~$M$) and its action is invisible on a single slot.
The fingerprint $\varphi$ is the minimal extension that cures
all three.
\end{remark}

\section{The fingerprint is a complete invariant}
\label{sec:fingerprint-is-complete-invariant}

\begin{lemma}[Presentation-independence of the fingerprint slots on the standard landscape]
\label{lem:fingerprint-slot-presentation-independence}
\ClaimStatusProvedHere
Let~$\cA$ be a chiral algebra in the standard landscape, and let
$\varphi(\cA)
= (p_{\max}, r_{\max}, \chi_{\mathrm{VOA}}, n_{\mathrm{strong}}, \mathrm{coset})$
be the canonical fingerprint of
Definition~\ref{def:canonical-fingerprint}. Then each of the five
slots is independent of the chosen minimal strong generating
presentation of~$\cA$.
\end{lemma}

\begin{proof}
Write
$\cA_+ := \bigoplus_{n>0}\cA_n$
for the positive-conformal-weight subspace and
\[
Q(\cA)
:=
\cA_+/\bigl(T\cA_+ + \cA_+(-1)\cA_+\bigr)
\]
for the indecomposable quotient. A minimal strong generating set of
standard-landscape~$\cA$ is precisely a basis of~$Q(\cA)$, so a change
of presentation is a basis change in~$Q(\cA)$.

\emph{Slot~(1): \(p_{\max}\).}
For each~$n\geq 1$, the OPE coefficient
$(a,b)\mapsto a_{(n-1)}b$ descends to a bilinear map
\[
\mu_n^\cA\colon Q(\cA)\otimes Q(\cA)\longrightarrow \gr\cA,
\]
because total derivatives and normally ordered composites are killed in
$Q(\cA)$. The condition that some generator pair has a non-zero
$1/(z-w)^n$-term is exactly the condition~$\mu_n^\cA\neq 0$, which is
basis-independent. Therefore
$p_{\max}(\cA)=\max\{n\mid \mu_n^\cA\neq 0\}$ is independent of the
chosen presentation; this sharpens the standard-landscape
well-definedness recorded in
Remark~\ref{rem:p-max-generator-choice}.

\emph{Slot~(2): \(r_{\max}\).}
By Theorem~\ref{thm:mc2-bar-intrinsic}, the full shadow obstruction
tower~$\Theta_\cA$ is the positive-genus correction of the bar
differential and is therefore intrinsic to~$\bar{B}^{\mathrm{ch}}(\cA)$.
Hence
\[
r_{\max}(\cA)=\max\{r\geq 2\mid \pi_r(\Theta_\cA)\neq 0\}
\]
is independent of any chosen presentation.

\emph{Slot~(3): \(\chi_{\mathrm{VOA}}\).}
The Hodge-filtered supertrace
$\chi_{\mathrm{VOA}}(\cA)=\operatorname{str}_{F^0\!\cA}(q^{L_0})$
depends only on the
$L_0$-graded super vector space underlying~$\cA$. Changing generators
does not change the eigenspace decomposition of~$L_0$ or the
$\mathbb{Z}/2$-parity, so the coefficients of~$\chi_{\mathrm{VOA}}$ are
presentation-independent.

\emph{Slot~(4): \(n_{\mathrm{strong}}\).}
Minimal strong generating sets are bases of~$Q(\cA)$, hence
\[
n_{\mathrm{strong}}(\cA)=\dim Q(\cA),
\]
which is basis-independent.

\emph{Slot~(5): \(\mathrm{coset}\).}
When~$\kappa_{\mathrm{ch}}(\cA)\neq 0$, the Heisenberg factor
$\cH_{\kappa_{\mathrm{ch}}}\hookrightarrow \cA$ is the intrinsic abelian
current subalgebra singled out by the invariant bilinear form on the
weight-one sector, and the pair
$(\mathrm{Com}(\cH_{\kappa_{\mathrm{ch}}},\cA),\,
\cA/\mathrm{Com}(\cH_{\kappa_{\mathrm{ch}}},\cA))$
is functorial in that inclusion. At~$\kappa_{\mathrm{ch}}(\cA)=0$,
Theorem~\ref{thm:fifth-class-FF}(iv) identifies the critical central
subalgebra with
$\operatorname{Fun}(\mathrm{Op}_{\fg^\vee}(D))$, so the critical coset
slot~$\mathrm{coset}_{FF}$ is likewise intrinsic. Thus the coset slot is
presentation-independent on the whole standard landscape.
\end{proof}

\begin{theorem}[Infinite fingerprint classification; strengthening of
Theorem~\ref{thm:fingerprint-completeness}]
\label{thm:fingerprint-is-complete-invariant}
\ClaimStatusProvedHere
Let~$\cA_1,\cA_2$ be chiral algebras in the standard landscape. The
following are equivalent:
\begin{enumerate}[label=\textup{(\alph*)},leftmargin=*]
\item the canonical fingerprints agree,
$\varphi(\cA_1) = \varphi(\cA_2)$;
\item the bar coalgebras are isomorphic,
$\bar{B}^{\mathrm{ch}}(\cA_1) \simeq \bar{B}^{\mathrm{ch}}(\cA_2)$
as $\mathrm{Ass}^{\mathrm{ch}}$-coalgebras, and the Koszul duals
are isomorphic,
$\cA_1^{\,!} \simeq \cA_2^{\,!}$
as chiral algebras with~$\mathrm{Com}^{\mathrm{ch}}$-structure
\textup{(}on the Koszul locus\textup{)};
\item the shadow obstruction towers agree,
$\Theta_{\cA_1}^{\leq r} \simeq \Theta_{\cA_2}^{\leq r}$ for
every~$r$, as
$(\mathrm{SC}^{\mathrm{ch,top}})^!$-coalgebras.
\end{enumerate}
\end{theorem}

\begin{proof}
\textit{Step 0: the five slots are intrinsic.}
Lemma~\ref{lem:fingerprint-slot-presentation-independence} removes the
presentation ambiguity of all five fingerprint slots on the
standard-landscape scope of the theorem.

\emph{Step 1: \textup{(a)} implies \textup{(b)}.}
Assume~$\varphi(\cA_1)=\varphi(\cA_2)$. By
Theorem~\ref{thm:fingerprint-completeness}, there is an
$\mathrm{Ass}^{\mathrm{ch}}$-coalgebra isomorphism
\[
f\colon \bar{B}^{\mathrm{ch}}(\cA_1)\xrightarrow{\ \sim\ }
\bar{B}^{\mathrm{ch}}(\cA_2)
\]
compatible with the
$(\mathrm{SC}^{\mathrm{ch,top}})^!$-coalgebra structures on the shadow
towers. Apply the chiral cobar functor~$\Omega$ to obtain
\[
\Omega(f)\colon
\Omega\bigl(\bar{B}^{\mathrm{ch}}(\cA_1)\bigr)\xrightarrow{\ \sim\ }
\Omega\bigl(\bar{B}^{\mathrm{ch}}(\cA_2)\bigr).
\]
On the Koszul locus, Chapter~\ref{ch:cobar} provides the bar--cobar
comparison isomorphisms
\[
\eta_i\colon
\Omega\bigl(\bar{B}^{\mathrm{ch}}(\cA_i)\bigr)\xrightarrow{\ \sim\ }\cA_i^{\,!},
\qquad i=1,2.
\]
Hence
\[
\eta_2\circ \Omega(f)\circ \eta_1^{-1}\colon
\cA_1^{\,!}\xrightarrow{\ \sim\ }\cA_2^{\,!}
\]
is the required Koszul-dual isomorphism. This proves~\textup{(a)}
$\Rightarrow$~\textup{(b)}.

\emph{Step 2: \textup{(b)} implies \textup{(c)}.}
The shadow tower is bar-intrinsic
\textup{(}Theorem~\ref{thm:mc2-bar-intrinsic}\textup{)}: for
every~$r\geq 2$,
\[
\Theta_\cA^{\leq r}
=
\pi_{\leq r}\bigl(\mathrm{MC}(\bar{B}^{\mathrm{ch}}(\cA))\bigr),
\]
where~$\pi_{\leq r}$ is the weight-$r$ truncation. The coalgebra
isomorphism in~\textup{(b)} transports Maurer--Cartan solutions
functorially, and the
$(\mathrm{SC}^{\mathrm{ch,top}})^!$-coaction comes from the ambient
operad, not from a chosen presentation of~$\cA$. Therefore
\textup{(b)}~$\Rightarrow$~\textup{(c)}.

\emph{Step 3: \textup{(c)} implies \textup{(a)}.}
Read the five slots from the intrinsic tower, with
$Q(\cA)=\cA_+/\bigl(T\cA_+ + \cA_+(-1)\cA_+\bigr)$ as in
Lemma~\ref{lem:fingerprint-slot-presentation-independence}.
\begin{enumerate}[label=\textup{(\arabic*)},leftmargin=*]
\item The OPE pole order~$p_{\max}$ is the maximal pole appearing in
the bilinear OPE maps
\[
\mu_n^\cA\colon Q(\cA)\otimes Q(\cA)\to \gr\cA,
\qquad
\mu_n^\cA([a],[b])=[a_{(n-1)}b],
\]
encoded by the degree-$2$ truncation~$\Theta_\cA^{\leq 2}$. Thus
$p_{\max}(\cA)=\max\{n\mid \mu_n^\cA\neq 0\}$ is preserved by an
isomorphism of shadow towers.
\item The shadow depth~$r_{\max}$ is the largest~$r$ for which
$\pi_r(\Theta_\cA)\neq 0$, so it is the termination depth of the tower
itself.
\item The Hodge-filtered supertrace~$\chi_{\mathrm{VOA}}$ is
the generating function
$\sum_n (-1)^{|n|}\dim(F^0\cA)_n q^n$, where the
$L_0$-graded super vector space~$F^0\cA$ is the degree-$1$ cogenerator
part of the tower. An isomorphism of towers preserves this graded
superdimension, hence preserves~$\chi_{\mathrm{VOA}}$.
\item The number~$n_{\mathrm{strong}}$ of minimal strong generators
is~$\dim Q(\cA)$, and the degree-$1$ cogenerator space of the tower is
exactly~$Q(\cA)$.
\item The coset slot~$\mathrm{coset}(\cA)$ is the isomorphism class
of the pair
$(\mathrm{Com}(\cH_{\kappa_{\mathrm{ch}}},\cA),\cA/\mathrm{Com})$
\textup{(}or~$\mathrm{coset}_{FF}$ at~$\kappa_{\mathrm{ch}}=0$\textup{)},
both of which are internal to the bar/tower datum: on the non-critical
locus the Heisenberg sub-pair is the image of the degree-$2$ center of
the tower under the natural projection, while at~$\kappa_{\mathrm{ch}}=0$
Theorem~\ref{thm:fifth-class-FF}(iv) identifies the same intrinsic
center with the oper algebra
$\operatorname{Fun}(\mathrm{Op}_{\fg^\vee}(D))$.
\end{enumerate}
Thus~\textup{(c)}~$\Rightarrow$~\textup{(a)}.

\emph{Step 4: close the biconditional.}
Combining Steps~2 and~3 gives
\textup{(b)}~$\Rightarrow$~\textup{(a)}$;$ the extra Koszul-dual
isomorphism in~\textup{(b)} is therefore compatible with, but not needed
for, the reverse implication. Together with Step~1 this supplies both
directions of the biconditional
\[
\varphi(\cA_1)=\varphi(\cA_2)
\Longleftrightarrow
\bigl(
\bar{B}^{\mathrm{ch}}(\cA_1)\simeq \bar{B}^{\mathrm{ch}}(\cA_2)
\text{ and, on the Koszul locus, }
\cA_1^{\,!}\simeq \cA_2^{\,!}
\bigr),
\]
and hence the equivalence of~\textup{(a)}, \textup{(b)}, and
\textup{(c)}.
\end{proof}

\section{The critical boundary stratum: critical-level bar structure}
\label{sec:FF-bar-structure}

The refined FF stratum is not the absence of a bar complex, but the
critical fixed-point geometry of the affine Kac--Moody family: the bar
complex of~$V_{-h^\vee}(\fg)$ is well-defined and uncurved, the center
jumps to the Feigin--Frenkel oper algebra, and the generic class-L
Theorem~H / raw bar-cobar package must be replaced by its critical
weight-completed form. What changes is not the existence of the bar
complex, but the surrounding theorem surface.

\begin{theorem}[Feigin--Frenkel stratum at the critical affine fixed point;
\ClaimStatusProvedHere]
\label{thm:fifth-class-FF}
Let $\fg$ be a simple Lie algebra and set
\[
\cA_{\mathrm{FF}} := V_{-h^\vee}(\fg).
\]
Then:
\begin{enumerate}[label=\textup{(\roman*)},leftmargin=*]
\item \textup{(}Position in the classification\textup{)}
 the coarse shadow-depth projection remains the affine
 Lie/tree value $r_{\max}=3$: the cubic shadow is the Lie bracket and
 the quartic shadow vanishes by the Jacobi identity. Hence FF is not a
 fifth \emph{generic} shadow-depth class; it is the critical
 fixed-point boundary stratum of class~L, recorded separately only in
 the refined fingerprint / atlas surface.
\item \textup{(}Scalar invariants\textup{)}
 \[
 \kappa(\cA_{\mathrm{FF}})
 = \frac{\dim(\fg)\,(-h^\vee+h^\vee)}{2h^\vee}
 = 0,
 \qquad
 \kappa(\cA_{\mathrm{FF}}^{!}) = 0,
 \]
 because the Feigin--Frenkel level involution
 $k \mapsto -k-2h^\vee$ fixes $k=-h^\vee$. Consequently
 \[
 \kappa(\cA_{\mathrm{FF}}) + \kappa(\cA_{\mathrm{FF}}^{!}) = 0,
 \qquad
 \operatorname{obs}_g(\cA_{\mathrm{FF}})
 = \kappa(\cA_{\mathrm{FF}})\lambda_g = 0
 \quad (g \ge 1).
 \]
\item \textup{(}Sugawara degeneration and bar-cobar scope\textup{)}
 the Sugawara field
 \[
 T_{\mathrm{Sug}}
 = \frac{1}{2(k+h^\vee)}\sum_a {:}J^aJ_a{:}
 \]
 is undefined at $k=-h^\vee$. Accordingly the raw non-critical
 class-L topologization and raw direct-sum Positselski / bar-cobar
 statements do not extend verbatim through the critical point; the
 surviving critical statement is the weight-completed / coderived
 bar-cobar package.
\item \textup{(}Center, derived center, and opers\textup{)}
 the ordinary center is the Feigin--Frenkel center
 \[
 Z(\cA_{\mathrm{FF}})
 \cong \operatorname{Fun}\!\bigl(\mathrm{Op}_{\fg^\vee}(D)\bigr)
 \]
 \textup{(}Feigin--Frenkel~\cite{Feigin-Frenkel}, Frenkel~\cite{Fre07},
 Beilinson--Drinfeld~\cite{BD04}\textup{)}.
 The critical bar complex computes the oper de~Rham package:
 \[
 H^n\!\bigl(\bar{B}^{\mathrm{ch}}(\cA_{\mathrm{FF}})\bigr)
 \cong \Omega^n\!\bigl(\mathrm{Op}_{\fg^\vee}(D)\bigr)
 \qquad (n \ge 0),
 \]
 and the derived center is computed in Vol~I via this same critical bar
 complex as a bar-Ext resolution
 \textup{(}Theorem~\textup{\ref{thm:langlands-bar-bridge}(iii)}\textup{)}.
\item \textup{(}Theorem~H boundary\textup{)}
 the generic Theorem~H concentration in~$\{0,1,2\}$ does not persist
 on FF: the chiral Hochschild theory at critical level is unbounded,
 with the degree-zero piece already infinite-dimensional through the
 Feigin--Frenkel center.
\item \textup{(}Geometric Langlands face\textup{)}
 the oper-algebra center is already present in Vol~I: the chapter on
 derived Langlands identifies the critical bar complex with the oper
 differential-form package and records this as the manuscript's
 internal geometric-Langlands bridge.
\end{enumerate}
\end{theorem}

\begin{proof}
\emph{Path~1: direct affine formulas.}
The affine Kac--Moody census gives
$\kappa(V_k(\fg)) = \dim(\fg)(k+h^\vee)/(2h^\vee)$, so
$\kappa(\cA_{\mathrm{FF}})=0$ at $k=-h^\vee$
\textup{(}Chapter~\ref{ch:landscape-census},
Chapter~\ref{chap:shadow-quadrichotomy-platonic}\textup{)}.
The Feigin--Frenkel involution $k \mapsto -k-2h^\vee$ fixes
$k=-h^\vee$, hence $\kappa(\cA_{\mathrm{FF}}^{!})=0$ as well and the
complementarity sum is~$0$. The genus-universality formula
$\operatorname{obs}_g = \kappa\lambda_g$ then gives
$\operatorname{obs}_g(\cA_{\mathrm{FF}})=0$ for all~$g \ge 1$.

\emph{Path~2: shadow/OPE analysis.}
For affine Kac--Moody the cubic shadow is the Lie bracket and the
quartic shadow vanishes by the Jacobi identity; the Kac--Moody census
and archetype table therefore keep the shadow-depth value
$r_{\max}=3$ at the critical fixed point. The same OPE analysis shows
that the Sugawara denominator $2(k+h^\vee)$ vanishes at $k=-h^\vee$,
so the non-critical topologization proof cannot continue through the
fixed point.

\emph{Path~3: critical-level bar / oper package.}
The Feigin--Frenkel theorem identifies
$Z(\cA_{\mathrm{FF}})$ with
$\operatorname{Fun}(\mathrm{Op}_{\fg^\vee}(D))$
\textup{(}\cite{Feigin-Frenkel}, \cite{Fre07}\textup{)}.
Theorem~\ref{thm:langlands-bar-bridge} upgrades this inside Vol~I:
the critical bar complex is uncurved and its cohomology is the oper
differential-form algebra in all degrees, while clause~(iii) records
the derived-center computation via the same critical bar complex.
Theorem~H's generic concentration is excluded at critical level, and
Remark~\ref{rem:critical-level-lie-vs-chirhoch} records the surviving
critical statement: chiral Hochschild cohomology is unbounded.
Finally, Corollary~\ref{cor:positselski-applicable-families} and
Remark~\ref{rem:ainf2-heal-74} give the correct Theorems~A/B scope at
the critical point: the raw direct-sum theorem is replaced by the
weight-completed / coderived package. This proves all clauses.
\end{proof}

\begin{remark}[Three consistency attacks on the FF stratum]
\label{rem:ff-consistency-attacks}
The healed FF characterization survives three independent attacks.
\begin{enumerate}[label=\textup{(\arabic*)},leftmargin=*]
\item \emph{$r_{\max}$ attack.}
 If FF were a new generic shadow-depth class, one would need a new
 quartic or higher witness. None appears: the affine cubic shadow
 survives and the Jacobi identity still kills the quartic term, so the
 coarse depth value remains~$3$.
\item \emph{Theorems~C/D/H attack.}
 The self-dual critical level gives $\kappa^!=0$ and therefore
 $\operatorname{obs}_g=0$, but the Feigin--Frenkel center forces
 unbounded chiral Hochschild cohomology. Thus FF is compatible with
 Theorem~D only after one excludes the generic Theorem~H surface and
 replaces it by the critical bar/oper package.
\item \emph{Theorems~A/B attack.}
 The Sugawara denominator shows that the raw non-critical class-L
 bar-cobar/topologization proof breaks at $k=-h^\vee$. What survives
 is exactly the weight-completed / coderived critical regime already
 inscribed elsewhere in Vol~I.
\end{enumerate}
No further structural gap remains on this theorem surface: every
surviving assertion is either proved above or explicitly scoped to the
completed critical regime.
\end{remark}

\begin{remark}[AP\textup{77} consequence: W(p) has $4$ strong
generators]
\label{rem:ap77-w-p-strong-gen}
Slot~(4) of the fingerprint records the \emph{strong}-generator
count, not the \emph{free-strong} count. For the logarithmic
$\Walg(p)$-algebra, the strong-generating set is $\{T, W^-, W^0, W^+\}$
(the Virasoro field plus an~$\mathfrak{sl}_2$-triplet of
weight-$(2p-1)$ primaries), giving $n_{\mathrm{strong}}(\Walg(p)) = 4$.
The free-strong question (whether these generate freely as a chiral
algebra) is distinct and open. The fingerprint uses the strong count
because this is the invariant accessible from the bar coalgebra
via rank of the degree-$1$ graded piece.
\end{remark}

\section{Pole--depth independence}
\label{sec:pole-depth-independence}

\begin{theorem}[Pole--depth independence]
\label{thm:pole-depth-independence}
\ClaimStatusProvedHere
The pair $(p_{\max}, r_{\max})$ takes the following values on the
standard landscape, which demonstrate that $p_{\max}$ and $r_{\max}$
vary independently:
\begin{center}
\renewcommand{\arraystretch}{1.3}
\begin{tabular}{lcccl}
Family & $p_{\max}$ & $r_{\max}$ & Class & Key OPE feature \\
\hline
Heisenberg~$\cH_k$ & $2$ & $2$ & $G$ & $JJ\sim k/(z-w)^2$ \\
$\beta\gamma$ & $1$ & $4$ & $C$ & $\beta\gamma\sim 1/(z-w)$ \\
$bc$ ($\lambda = 2$) & $1$ & $4$ & $C$ & $bc\sim 1/(z-w)$ \\
Affine $\widehat{\fg}_k$ & $2$ & $3$ & $L$ & $JJ$ current \\
Virasoro $\mathrm{Vir}_c$ & $4$ & $\infty$ & $M$ & $TT$ stress \\
$\Walg_3$ & $6$ & $\infty$ & $M$ & $WW$ cubic \\
$\Walg_N$ & $2N$ & $\infty$ & $M$ & $W_NW_N$ \\
	$V_{-h^\vee}(\fg)$ & $2$ & $3$ & FF boundary of $L$ & Sugawara undefined
	\end{tabular}
	\end{center}
	In particular, the pairs $(2,2)$ and $(2,3)$ show that $p_{\max}=2$
	does not determine~$r_{\max}$; the pairs $(1,4)$ and $(4,\infty)$
	show that the inequality~$p_{\max} \lessgtr r_{\max}$ can go either
	way. The FF row records a different phenomenon: the critical affine
	fixed point keeps the coarse Lie/tree value $r_{\max}=3$ while the
	Sugawara topologization degenerates, so FF is a critical boundary
	stratum rather than a new generic $r_{\max}=\infty$ family.
\end{theorem}

\begin{proof}
Each row is a citation to the landscape census
	\textup{(}Chapter~\ref{ch:landscape-census}\textup{)}: Heisenberg
	\eqref{eq:heisenberg-kappa}, $\beta\gamma$
	\eqref{eq:beta-gamma-ope}--\eqref{eq:beta-gamma-r-max}, affine
	\eqref{eq:affine-KM-p-max}, Virasoro \eqref{eq:vir-p-max},
	$\Walg_N$ \eqref{eq:w-N-p-max}, critical affine via
	Theorem~\ref{thm:fifth-class-FF}\textup{(}i,iii\textup{)}. Independence is the
	existence of the table.
\end{proof}

\begin{proposition}[Class determination from full shadow tower;
formalisation of AP-CY\textup{12}]
\label{prop:class-from-full-tower}
\ClaimStatusProvedHere
The shadow depth~$r_{\max}(\cA)$ is computed from the full shadow
tower~$\Theta_\cA^{\leq r}$ for~$r = 2, 3, 4, 5$ and the
non-termination witness at~$r = 5$; it is not determined by~$p_{\max}$
alone, by~$k_{\max}$ alone, or by counting strong generators. The
$G/L/C/M$ classification therefore requires the full MC-tower
computation, consistent with AP-CY\textup{12}.
\end{proposition}

\begin{proof}
Theorem~\ref{thm:pole-depth-independence} exhibits pairs with
equal~$p_{\max}$ and distinct~$r_{\max}$ (Heisenberg versus affine
at~$p_{\max} = 2$), and pairs with equal~$k_{\max}$ and distinct
$r_{\max}$. Counting strong generators fails because~$\Walg_3$ has
more strong generators than~$\mathrm{Vir}_c$ yet both lie in
class~$M$; and~$\Walg(p)$ has four strong generators yet its
class is controlled by the shadow tower, not the count
\textup{(}Remark~\ref{rem:ap77-w-p-strong-gen}\textup{)}. The
non-termination witness at~$r = 5$ follows from the recursive
MC-inverse-limit construction~(Theorem~\ref{thm:mc2-bar-intrinsic})
evaluated at degree~$5$: the family has~$r_{\max}<\infty$ iff the
degree-$5$ component of~$\Theta_\cA$ is annihilated by the
Maurer--Cartan closure at that degree.
\end{proof}

\section{The depth-bijection: $d_{\mathrm{alg}} = r_{\max} - 2$}
\label{sec:d-alg-r-max-bijection}

\begin{theorem}[Depth bijection; closes FM\textup{110}]
\label{thm:d-alg-r-max-bijection}
\ClaimStatusProvedHere
For every chirally Koszul algebra~$\cA$ with $\kappa_{\mathrm{ch}}(\cA)\neq 0$,
\begin{equation}\label{eq:d-alg-r-max}
d_{\mathrm{alg}}(\cA) \;=\; r_{\max}(\cA) - 2
\quad\in\quad \{0, 1, 2, \infty\}.
\end{equation}
In particular, the depth gap $d_{\mathrm{alg}}\in\{0,1,2,\infty\}$
of Proposition~\ref{prop:depth-gap-trichotomy} is the
quadrichotomy~$\{G,L,C,M\}$ shifted by~$-2$.
\end{theorem}

\begin{proof}
Recall $d_{\mathrm{alg}}(\cA)$ is the highest weight at which the
MC tower~$\Theta_\cA$ carries a nonzero obstruction
\textup{(}Definition~\ref{def:d-alg}\textup{)}, while $r_{\max}(\cA)$
is the highest bar-weight at which the shadow obstruction
$\Theta_\cA^{\leq r}$ differs from~$\Theta_\cA^{\leq r-1}$.

The bar complex has an internal weight shift of~$+2$: the primitive
obstruction at algebraic weight~$n$ enters the bar complex at
bar-weight~$n+2$, because the first two bar-degrees are
reserved for the Maurer--Cartan equation itself
$(D^*\Theta + \tfrac{1}{2}[\Theta,\Theta] = 0)$ and the
coboundary compatibility. Explicitly, for each of the four classes:
\begin{enumerate}[label=\textup{(\alph*)},leftmargin=*]
\item Class~$G$: $d_{\mathrm{alg}} = 0$, $r_{\max} = 2$. The
Heisenberg shadow tower terminates at bar weight~$2$ with
$\Theta^{\leq 2} = \kappa_{\mathrm{ch}}$; no higher algebraic
obstruction.
\item Class~$L$: $d_{\mathrm{alg}} = 1$, $r_{\max} = 3$. Affine
Kac--Moody has a weight-$1$ Lie obstruction
\textup{(}the structure constants~$f^{abc}$\textup{)}, which enters
the bar at weight~$3$ via the Arnold relation.
\item Class~$C$: $d_{\mathrm{alg}} = 2$, $r_{\max} = 4$. The
$\beta\gamma$ system has a weight-$2$ quartic contact obstruction,
entering the bar at weight~$4$.
\item Class~$M$: $d_{\mathrm{alg}} = \infty$, $r_{\max} = \infty$.
The Virasoro shadow tower is infinite; the depth gap
\textup{(}Proposition~\ref{prop:depth-gap-trichotomy}\textup{)}
shows~$d_{\mathrm{alg}} = 3$ is impossible, so
$d_{\mathrm{alg}}\in\{0,1,2\}\cup\{\infty\}$.
\end{enumerate}
The~$+2$ shift is uniform across all four classes and across the
standard landscape, giving~\eqref{eq:d-alg-r-max}. At~$\kappa_{\mathrm{ch}} = 0$
\textup{(}class~$\mathrm{FF}$\textup{)} the algebraic depth is
ill-defined because the Sugawara construction degenerates, which is
why~\eqref{eq:d-alg-r-max} is stated on~$\kappa_{\mathrm{ch}}\neq 0$.
\end{proof}

\begin{corollary}[Quadrichotomy as depth shift]
\label{cor:quadrichotomy-depth-shift}
The $G/L/C/M$ classification by $r_{\max}\in\{2,3,4,\infty\}$ is
equivalent to classification by
$d_{\mathrm{alg}}\in\{0,1,2,\infty\}$ on the non-critical locus.
\end{corollary}

\section{The coarse projection and its non-trivial fibres}
\label{sec:coarse-projection-fibres}

\begin{theorem}[Quadrichotomy is a coarse projection; strengthening
of Proposition~\ref{prop:coarse-projection-functor}]
\label{thm:quadrichotomy-is-coarse-projection}
\ClaimStatusProvedHere
The coarse projection
\begin{equation}\label{eq:coarse-projection-fibres}
\Pi_{\mathrm{coarse}} \colon \mathcal{F} \twoheadrightarrow \{G, L, C, M, \mathrm{FF}\}
\end{equation}
is surjective with non-trivial fibres: each of the five classes
contains at least two fingerprints distinguished by slots~$(1),(3),(5)$.
Explicit witness pairs are:
\begin{enumerate}[label=\textup{(\roman*)},leftmargin=*]
\item \emph{Class~$G$ fibre:} Heisenberg~$\cH_k$ and abelian
$\widehat{\fu(1)}_{k'}$ share $(r_{\max} = 2, \chi_{\mathrm{VOA}})$
up to sign; they differ in $p_{\max}$ (both~$2$ but at different
levels) and in the coset slot
\textup{(}$\mathrm{Sp}(2,\mathbb{R})$-centraliser versus abelian\textup{)}.
\item \emph{Class~$G$ fibre via parity:} symplectic boson~$\beta\gamma_{1/2}$
and symplectic fermion~$\mathcal{SF}$ share the same coarse class
under shadow depth~$r_{\max} = 2$ at the weight-$1/2$ primary,
yet their fingerprints differ in sign of~$\chi_{\mathrm{VOA}}$
and in the coset pair~$\mathrm{Sp}(2)$ versus~$\mathrm{OSp}(1|2)$.
\item \emph{Class~$C$ fibre:} the bosonic quartic-contact algebra
$\Walg_4^{\mathrm{bos}}$ and the super-$\Walg_3$ algebra share
$r_{\max} = 4$, $p_{\max} = 8$, $n_{\mathrm{strong}} = 4$; they
differ in $\chi_{\mathrm{VOA}}$ (the super variant has an odd
weight-$3/2$ generator contributing negatively).
\end{enumerate}
Consequently, the quadrichotomy $\Pi_{\mathrm{coarse}}\circ\varphi$
is strictly lossy.
\end{theorem}

\begin{proof}
Surjectivity is the landscape census. Non-trivial fibres are the
three witness pairs. Each pair is verified by reading the five
fingerprint slots: the shadow-depth slot~$(2)$ agrees, one or more
of slots~$(1),(3),(5)$ disagrees.

\emph{Witness (i): Heisenberg versus abelian affine.}
$\varphi(\cH_k) = (2, 2, \operatorname{ch}_q\cH_k, 1, \mathrm{Sp}(2,\mathbb{R}))$
versus
$\varphi(\widehat{\fu(1)}_{k'}) = (2, 2, \operatorname{ch}_q\widehat{\fu(1)}, 1, \mathrm{U}(1))$.
Slot~(5) distinguishes: $\widehat{\fu(1)}$ has the additive group
$\mathrm{U}(1)$ as coset, whereas Heisenberg carries the full
symplectic structure.

\emph{Witness (ii): symplectic boson versus symplectic fermion.}
Both have $(p_{\max}, r_{\max}) = (1, 4)$ and live in class~$C$
under the full shadow tower, not class~$G$
\textup{(}Example~\ref{ex:symplectic-boson-vs-fermion}\textup{)}.
The parity slot separates:
$\chi_{\mathrm{VOA}}(\beta\gamma_{1/2})$ has a positive weight-$1/2$
boson, $\chi_{\mathrm{VOA}}(\mathcal{SF})$ has a negative weight-$1$
fermion. The coset slot also separates.

\emph{Witness (iii): bosonic $\Walg_4^{\mathrm{bos}}$ versus super-$\Walg_3$.}
Both have shadow depth~$4$ via the~$\mathrm{Vir}\otimes\Walg_3$-type
quartic obstruction, both have four strong generators (stress
tensor plus three primaries), yet the super variant includes a
weight-$3/2$ fermionic primary which contributes
$\operatorname{str}(q^{L_0}) \ni -q^{3/2}\mathbb{Z}[\![q]\!]_+$.
This forces $\chi_{\mathrm{VOA}}$ to differ.
\end{proof}

\begin{remark}[Witness (iii) closes FM\textup{107}]
\label{rem:fm107-closure}
FM\textup{107} warned that depth-$4$ bosonic $W$-algebras could be
conflated with super-$\Walg_3$ under the quadrichotomy. Witness
(iii) of Theorem~\ref{thm:quadrichotomy-is-coarse-projection}
produces the separation at the fingerprint level: slot~(3)
($\chi_{\mathrm{VOA}}$) and slot~(5) (coset) carry the
$\Zmod 2$-grading information that the quadrichotomy discards.
\end{remark}

\begin{remark}[Witness (ii) closes FM\textup{106}]
\label{rem:fm106-closure}
FM\textup{106} warned that the symplectic boson and the symplectic
fermion could be confused when reading shadow depth alone. The two
lie in class~$C$ (not class~$G$) under the full shadow tower;
within class~$C$ they are separated by slots~(3) and~(5).
\end{remark}

\section{DS transport of the fingerprint}
\label{sec:DS-transport}

Drinfeld--Sokolov reduction sends affine Kac--Moody algebras to
$W$-algebras. On fingerprints, this is an explicit slot-by-slot
transport: the coset slot transforms predictably, the pole order
changes in accordance with the new stress tensor, the shadow depth
escalates from $3$ to $\infty$ (class $L$ to class $M$), and the
character transforms as a BRST supertrace. Two confusions are cleared
up: the shadow depth \emph{must} escalate because the stress tensor
is a polynomial in $J^a$, not a primary field of~$\widehat\fg$; and
the strong-generator count grows by the number of simple roots
corresponding to the principal nilpotent (or to the chosen nilpotent
orbit in the non-principal case).

\begin{theorem}[DS transport of the fingerprint; closes FM\textup{108}]
\label{thm:DS-fingerprint-transport}
\ClaimStatusProvedHere
Let $\fg$ be a simple Lie algebra of rank~$r$ with Coxeter
number~$h^\vee$, and let~$\Gamma \subset \fg$ be the principal
nilpotent. The Drinfeld--Sokolov reduction
$H^\bullet_{\mathrm{DS},\Gamma}(\widehat{\fg}_k) = \Walg^\Gamma_k(\fg)$
induces, on fingerprints, the following slot-by-slot transport
\textup{(}at a generic non-critical level~$k\neq -h^\vee$\textup{)}:
\begin{enumerate}[label=\textup{(\roman*)},leftmargin=*]
\item \emph{Slot 1 (pole order):} $p_{\max}(\widehat{\fg}_k) = 2$ sends
to $p_{\max}(\Walg^\Gamma_k(\fg)) = 2h^\vee$ for~$\Gamma$ principal
in type~$A_{N-1}$: the stress-tensor OPE has pole order~$2\cdot 2 = 4$
and the highest-spin generator~$W^{(h^\vee)}$ has a self-OPE with
pole~$2h^\vee$.
\item \emph{Slot 2 (shadow depth):} $r_{\max}\colon 3 \mapsto \infty$.
Class $L$ escalates to class $M$; the escalation is strict because the
stress tensor of~$\Walg$ is a quadratic-in-$J$ polynomial, and its
self-OPE generates shadows at every weight.
\item \emph{Slot 3 (character):} $\chi_{\mathrm{VOA}}(\widehat{\fg}_k)$
sends to the BRST supertrace
$\chi_{\mathrm{VOA}}(\Walg^\Gamma_k(\fg)) = \operatorname{str}_{C^\bullet_{\mathrm{DS}}}(q^{L_0})$
where~$L_0$ is the DS-shifted grading by the principal
$\mathfrak{sl}_2$-element. Explicitly
$\chi \mapsto \chi_{\mathrm{DS}} = \chi\cdot\eta_{\mathrm{ghost}}$
where~$\eta_{\mathrm{ghost}}(q)$ is the Feigin--Frenkel ghost
character.
\item \emph{Slot 4 (strong-generator count):}
$n_{\mathrm{strong}}(\widehat{\fg}_k) = \dim\fg$ sends to
$n_{\mathrm{strong}}(\Walg^\Gamma_k(\fg)) = \mathrm{rank}(\fg) + \delta_\Gamma$,
with~$\delta_\Gamma = 0$ for principal~$\Gamma$ and $\delta_\Gamma \geq 0$
for non-principal nilpotent orbits; in type~$A_{N-1}$ with
$\Gamma$ principal, $n_{\mathrm{strong}} = N-1$.
\item \emph{Slot 5 (coset):} the affine coset
$(\mathrm{Com}(\cH, \widehat{\fg}_k), \widehat{\fg}_k/\mathrm{Com})$
transforms to the W-algebra coset
$(\mathrm{Com}(\cH, \Walg^\Gamma_k(\fg)), \Walg^\Gamma_k/\mathrm{Com})$;
on the Koszul dual side this is Feigin--Frenkel duality
$\Walg^\Gamma_k(\fg) \leftrightarrow \Walg^{\Gamma^\vee}_{k'}(\fg^\vee)$
with~$(k + h^\vee)(k' + h^{\vee,\vee}) = 1$.
\end{enumerate}
In particular, DS reduction is a well-defined map
$\mathrm{DS}_\Gamma\colon \mathcal{F}_L \to \mathcal{F}_M$ on
fingerprint strata, and it is not surjective onto class~$M$ (only
the $W$-algebras lying in the image of DS are reached; exceptional
$\Walg(p)$-type algebras are not).
\end{theorem}

\begin{proof}
Slot~(1) is the standard calculation for~$\Walg_N$: the highest-spin
generator has conformal weight~$N$, so the self-OPE has pole
order~$2N = 2h^\vee$ in type~$A_{N-1}$
\textup{(}\cite{BouwknegtSchoutens1993}\textup{)}.

Slot~(2) is the shadow-tower non-termination: $\Walg_N$ lies in class~$M$
for all~$N\geq 2$ and generic~$c$
\textup{(}Chapter~\ref{ch:w-algebras}\textup{)}. The escalation from
class~$L$ is strict because the Sugawara stress tensor of~$\widehat{\fg}_k$
is a \emph{secondary} field for the purposes of shadow classification
(it is not among the generators), whereas the stress tensor
of~$\Walg^\Gamma_k(\fg)$ is a primary strong generator.

Slot~(3) follows from the BRST reduction:
$C^\bullet_{\mathrm{DS}} = \widehat{\fg}_k \otimes \Lambda^\bullet(\fn_+^* \oplus \fn_+)$
\textup{(}\cite{deBoerTjin1993},~\cite{FrenkelBenZvi2004}\textup{)}
and the graded supertrace factorises
$\operatorname{str}(q^{L_0}) = \operatorname{ch}_{\widehat{\fg}_k}(q)\cdot\eta_{\mathrm{ghost}}(q)$
on the character level, then passes to cohomology on the BRST level.

Slot~(4) is Kac--Roan--Wakimoto \cite{KacRoanWakimoto2003}: the
principal DS reduction in type~$A_{N-1}$ produces $\Walg_N$ with
exactly~$N-1$ strong generators $W^{(2)},\ldots,W^{(N)}$; more
generally, the strong-generator count for non-principal~$\Gamma$
equals~$\mathrm{rank}(\fg) + $ a correction~$\delta_\Gamma\geq 0$
tracking the non-principal sector
\textup{(}\cite{KacWakimoto2004}\textup{)}.

Slot~(5) is Feigin--Frenkel duality for $W$-algebras
\textup{(}\cite{FeiginFrenkel1992},~\cite{Arakawa2007}\textup{)}:
the Koszul-dual pair of~$\Walg^\Gamma_k(\fg)$ under the
$k \leftrightarrow k'$ involution is~$\Walg^{\Gamma^\vee}_{k'}(\fg^\vee)$,
which records the Langlands-dual group and the dual nilpotent orbit.
The coset slot transforms accordingly.

Non-surjectivity onto class~$M$: the triplet
$\Walg(p)$-algebras
(logarithmic, with four strong generators including an
$\mathfrak{sl}_2$-triplet of primaries) are not Drinfeld--Sokolov
reductions of any affine algebra; they are in class~$M$ but not in
the image of~$\mathrm{DS}$.
\end{proof}

\begin{remark}[DS transport at~$\kappa_{\mathrm{ch}} = 0$]
\label{rem:DS-at-critical}
At~$k = -h^\vee$, the DS reduction sends the critical affine fixed point
$V_{-h^\vee}(\fg)$, i.e.\ the FF boundary point of class~$L$, to
$\Walg^\Gamma_{-h^\vee}(\fg)$. The target keeps the coarse
$W$-family shadow-depth behaviour~$r_{\max}=\infty$, but its coset slot
degenerates to oper data via the Feigin--Frenkel duality
$\Walg^\Gamma_{-h^\vee}(\fg) \simeq \mathrm{Op}_{\fg^\vee}$: classical
opers for the Langlands-dual group. The fingerprint slot~(5)
therefore transports~$\mathfrak{z}(\widehat{\fg}) \leftrightarrow \Fun\mathrm{Op}_{\fg^\vee}(D)$
under DS, coherent with
Theorem~\ref{thm:fifth-class-FF}(ii). This is a critical-surface
statement about slot~(5), not a
promotion of FF to a fifth generic shadow-depth class.
\end{remark}

\section{Synthesis and witness table for the generic classes and critical surface}
\label{sec:five-class-witness}

The strengthening theorems of this chapter assemble into a complete
fingerprint witness table for the four generic shadow-depth classes
and the critical surface. Each row is a five-slot reading, verified
against the landscape census.

\begin{table}[h]
\centering
\renewcommand{\arraystretch}{1.3}
\begin{tabular}{@{}lccccl@{}}
Family & $p_{\max}$ & $r_{\max}$ & Parity & $n_{\mathrm{strong}}$ & Coset (abbrev.) \\
\hline
$\cH_k$ (Heis) & $2$ & $2$ & $+$ & $1$ & $\mathrm{Sp}(2,\mathbb{R})$ \\
$\widehat{\fu(1)}_{k'}$ & $2$ & $2$ & $+$ & $1$ & $\mathrm{U}(1)$ \\
$V_\Lambda$ (lattice) & $2h_\Lambda$ & $2$ & $+$ & $\mathrm{rank}(\Lambda)$ & $\Lambda^*/\Lambda$ \\
\hline
$\widehat{\mathfrak{sl}_2}_k$ & $2$ & $3$ & $+$ & $3$ & $\fg/\cH$ \\
$\widehat{\fg}_k$ & $2$ & $3$ & $+$ & $\dim\fg$ & $\fg/\cH$ \\
\hline
$\beta\gamma$ & $1$ & $4$ & $+$ & $2$ & $\mathrm{Sp}(2)$ \\
$bc$ ($\lambda = 2$) & $1$ & $4$ & $-$ & $2$ & $\mathrm{OSp}(1|2)$ \\
$\mathcal{SF}$ (symp.\ fermion) & $1$ & $4$ & $-$ & $2$ & $\mathrm{OSp}(1|2)$ \\
\hline
$\mathrm{Vir}_c$ & $4$ & $\infty$ & $+$ & $1$ & $\mathrm{Vir}_c^!$ \\
$\Walg_3$ & $6$ & $\infty$ & $+$ & $2$ & $\Walg_3^!$ \\
$\Walg_N$ & $2N$ & $\infty$ & $+$ & $N-1$ & $\Walg_N^!$ \\
super-$\Walg_3$ & $6$ & $\infty$ & $\pm$ & $2$ & super-$\Walg_3^!$ \\
$\Walg(p)$ (triplet) & $2(2p-1)$ & $\infty$ & $+$ & $4$ & $\Walg(p)^!$ \\
\hline
	$V_{-h^\vee}(\mathfrak{sl}_2)$ & $2$ & $3$ & $+$ & $3$ & $\Fun\mathrm{Op}_{\mathfrak{sl}_2}(D)$ \\
	$V_{-h^\vee}(\fg)$ & $2$ & $3$ & $+$ & $\dim\fg$ & $\Fun\mathrm{Op}_{\fg^\vee}(D)$ \\
	critical $\Walg_N$ & $2N$ & $\infty$ & $+$ & $N-1$ & critical oper-coset \\
	\end{tabular}
	\caption{Fingerprint witness table for the generic classes and the
	critical surface. Parity records the
	sign of the~$\chi_{\mathrm{VOA}}$ leading term; coset is abbreviated
	to the dominant datum (centraliser group on the Koszul locus;
	$\mathrm{Op}$-data on the critical surface).}
	\label{tab:fingerprint-stratum}
\end{table}

\begin{corollary}[Completeness on the standard landscape]
\label{cor:fingerprint-separates-landscape}
\ClaimStatusProvedHere
Table~\ref{tab:fingerprint-stratum} exhibits pairwise distinct
fingerprints for every pair of standard-landscape chiral algebras
listed, including pairs within a single coarse class:
\begin{itemize}[nosep]
\item Heisenberg versus $\widehat{\fu(1)}_{k'}$ (class~$G$,
	slot~$5$);
\item $\beta\gamma$ versus symplectic fermion (class~$C$, slots~$3,5$);
\item $\Walg_3$ versus super-$\Walg_3$ (class~$M$, slot~$3$);
\item $V_{-h^\vee}(\fg)$ versus critical $\Walg_N$
	(critical surface, slots~$1,4,5$).
\end{itemize}
Together with the stratum theorem
\textup{(}Theorem~\ref{thm:five-class-stratum}\textup{)}, this
establishes that~$\varphi$ is the canonical complete invariant of
a chiral algebra in the standard landscape.
\end{corollary}

\begin{proof}
Direct inspection of Table~\ref{tab:fingerprint-stratum}. Each row
is witnessed by the corresponding census family and evaluated
against the five-slot specification
\textup{(}Definition~\ref{def:canonical-fingerprint}\textup{)}.
\end{proof}

\subsection{The holomorphic \texorpdfstring{$c=24$}{c=24} census is a
structured-subset theorem}
\label{subsec:schellekens-structured-subset}

\begin{theorem}[Structured-subset derivation of the holomorphic
\texorpdfstring{$c=24$}{c=24} census; closes AP\textup{290}]
\label{thm:schellekens-structured-subset}
\ClaimStatusProvedHere
Let $\mathfrak{S}_{24}^{\mathrm{hol}}$ be the closed ambient of the
$71$ strongly rational holomorphic VOAs of central charge~$24$ on
Schellekens's list, viewed as chiral algebras~$A$ with fingerprint
$\varphi(A)$. Then:
\begin{enumerate}[label=\textup{(\roman*)},leftmargin=*]
\item \textup{(}Structured subsets\textup{)}
 there is a disjoint decomposition
 \[
 \mathfrak{S}_{24}^{\mathrm{hol}}
 =
 \mathfrak{S}_{24}^{\mathrm{Niem}}
 \;\sqcup\;
 \{V^\natural\}
 \;\sqcup\;
 \mathfrak{S}_{24}^{\mathrm{nonlat}},
 \]
 where, in Schellekens's Table~2 \textup{(}1993, §3, pp.~177--182; the
 numbered list on pp.~178--182\textup{)},
 \[
 \begin{aligned}
 \mathfrak{S}_{24}^{\mathrm{Niem}}
 = \{&
 1{:}\,\fu(1)^{24},\;
 15{:}\,A_{1,1}^{24},\;
 24{:}\,A_{2,1}^{12},\;
 30{:}\,A_{3,1}^{8},\;
 37{:}\,A_{4,1}^{6},\;
 42{:}\,D_{4,1}^{6},\\
 &
 43{:}\,A_{5,1}^{4}D_{4,1},\;
 46{:}\,A_{6,1}^{4},\;
 49{:}\,A_{7,1}^{2}D_{5,1}^{2},\;
 51{:}\,A_{8,1}^{3},\;
 54{:}\,D_{6,1}^{4},\\
 &
 55{:}\,A_{9,1}^{2}D_{6,1},\;
 58{:}\,E_{6,1}^{4},\;
 59{:}\,A_{11,1}D_{7,1}E_{6,1},\;
 60{:}\,A_{12,1}^{2},\;
 61{:}\,D_{8,1}^{3},\\
 &
 63{:}\,A_{15,1}D_{9,1},\;
 64{:}\,D_{10,1}E_{7,1}^{2},\;
 65{:}\,A_{17,1}E_{7,1},\;
 66{:}\,D_{12,1}^{2},\;
 67{:}\,A_{24,1},\\
 &
 68{:}\,E_{8,1}^{3},\;
 69{:}\,D_{16,1}E_{8,1},\;
 70{:}\,D_{24,1}\},
 \end{aligned}
 \]
 row~$0$ is the Monster slot $V_1=0$, and the non-lattice summand is
 the explicit complement
 \[
 \mathfrak{S}_{24}^{\mathrm{nonlat}}
 =
 \{2\text{--}14,\;16\text{--}23,\;25\text{--}29,\;31\text{--}36,\;
 38\text{--}41,\;44\text{--}45,\;47\text{--}48,\;50,\;52\text{--}53,\;
 56\text{--}57,\;62\}.
 \]
	\item \textup{(}Disjointness and exhaustiveness\textup{)}
	 no algebra in~$\mathfrak{S}_{24}^{\mathrm{hol}}$ lies in two of the
	 three subsets: the Monster singleton is characterized by~$V_1=0$,
	 the Niemeier rows are characterized by
	 $V_1\cong \fu(1)^{24}$ or by level-$1$ simply-laced ADE datum of total
	 rank~$24$, and every remaining Schellekens row is, by definition, in
	 $\mathfrak{S}_{24}^{\mathrm{nonlat}}$. Conversely, every algebra
	 in~$\mathfrak{S}_{24}^{\mathrm{hol}}$ lies in exactly one of these
	 subsets because the displayed row sets partition~$\{0,1,\dots,70\}$.
\item \textup{(}Counting corollary, not theorem\textup{)}
 the identity
 \[
 71 = 24 + 1 + 46
 \]
 is a downstream count of the structured subsets in~\textup{(i)};
 it is not itself the classification theorem.
\item \textup{(}Niemeier rank=count, three independent paths\textup{)}
 the equality
 \[
 24 = \operatorname{rank}(N)
 = \#\mathfrak{S}_{24}^{\mathrm{Niem}}
 \]
 is verified in three independent ways:
 \begin{enumerate}[label=\textup{(\alph*)},leftmargin=*]
 \item Niemeier's class-number computation gives exactly~$24$
  isomorphism classes of positive-definite even unimodular
  lattices of rank~$24$
  \textup{(}Niemeier 1973, J. Number Theory~5, p.~142, abstract; see
  also Conway--Sloane, Table~16.1, p.~407\textup{)};
 \item the displayed ADE list is an in-place lattice computation:
  there are exactly~$23$ rootful simply-laced level-$1$ rank-$24$
  rows with common-Coxeter decomposition, and the Leech row
  $1{:}\,\fu(1)^{24}$ is the unique rootless case;
 \item the Schellekens-table intersection above has exactly~$24$ rows:
  the abelian row~$1$ together with the~$23$ level-$1$ simply-laced
  Niemeier root-system rows.
 \end{enumerate}
\end{enumerate}
\end{theorem}

\begin{proof}
\emph{Step~1: identify the rows.}
Schellekens's classification is by the weight-one Lie algebra~$V_1$,
with the Monster row singled out by~$V_1=0$ and the remaining rows
listed in Table~2. The rows written in
\textup{\ref{thm:schellekens-structured-subset}(i)} are exactly the
abelian row and the level-$1$ simply-laced ADE rows whose total rank is
$24$. They are therefore the Niemeier-lattice candidates. The
complementary row set is listed explicitly, so no unnamed remainder is
left. These are precisely the three membership tests stated in
Theorem~\ref{thm:schellekens-structured-subset}\textup{(}ii\textup{)}:
\[
V_1=0,\qquad
V_1\cong \fu(1)^{24}\ \text{or level-$1$ simply-laced ADE of rank~$24$},
\qquad
\text{complementary Schellekens row}.
\]

\emph{Step~2: lattice identification of the $24$-summand.}
For every positive-definite even unimodular rank-$24$ lattice~$N$, the
Frenkel--Kac--Segal / FLM construction gives a holomorphic lattice VOA
$V_N$ of central charge~$24$ whose weight-one Lie algebra is the
corresponding level-$1$ simply-laced root system, or $\fu(1)^{24}$ in
the Leech case. Niemeier's class-number computation
\textup{(}J. Number Theory~5 \textup{(}1973\textup{)}, p.~142,
abstract\textup{)} and Conway--Sloane's Table~16.1 \textup{(}p.~407\textup{)}
supply exactly~$24$ such lattices. Conversely, every row in
$\mathfrak{S}_{24}^{\mathrm{Niem}}$ is one of those lattice VOAs,
because its level-$1$ ADE datum is one of the~$23$ explicitly listed
rank-$24$ root systems or the abelian Leech row.

\emph{Step~3: isolate the Monster singleton.}
Schellekens isolates the unique row with~$V_1=0$; FLM identify that row
with the Moonshine module~$V^\natural$
\textup{(}FLM88, Chapters~8--11; compare also M\"oller--Scheithauer
2023, Theorem~6.9, pp.~40--41\textup{)}. This gives the singleton
summand~$\{V^\natural\}$.

\emph{Step~4: identify the $46$-summand.}
What remains after removing the~$24$ lattice rows and the Monster row
is the explicit complement set written above, so its cardinality is
$71-24-1=46$. M\"oller--Scheithauer 2023 give an independent
generalised-deep-hole realization of this complement:
Theorem~6.6 \textup{(}p.~33\textup{)} identifies nonzero
$V_1$-rows with positive-rank deep holes, the abstract
\textup{(}p.~2\textup{)} and Theorem~6.7 \textup{(}pp.~34--35\textup{)}
construct the nonzero Lie structures uniformly, and
Theorem~6.9 \textup{(}pp.~40--41\textup{)} isolates the
$V_1=0$ branch whose orbifolds are all
$V^\natural$. Intersecting those results with the explicit
Niemeier row set above leaves exactly the~$46$ genuinely non-lattice
holomorphic $c=24$ cases.

\emph{Step~5: disjointness and exhaustiveness on the structured subset.}
The three membership tests from Step~1 are mutually exclusive.
The condition~$V_1=0$ singles out~$V^\natural$, so the Monster row
cannot lie in the Niemeier or non-lattice summands. Every element of
$\mathfrak{S}_{24}^{\mathrm{Niem}}$ has
$V_1\cong \fu(1)^{24}$ or level-$1$ simply-laced ADE of total
rank~$24$, hence cannot be the Monster row and cannot lie in the
explicit complementary row set defining
$\mathfrak{S}_{24}^{\mathrm{nonlat}}$. Finally,
$\mathfrak{S}_{24}^{\mathrm{nonlat}}$ is defined as that complementary
row set, so it contains none of the Niemeier or Monster rows.
Conversely, Schellekens's table lists all rows~$0,\dots,70$, and the
three displayed row sets cover them all. Therefore every algebra in
$\mathfrak{S}_{24}^{\mathrm{hol}}$ lies in exactly one of the three
structured subsets.

\emph{Step~6: rank=count.}
Path~(a) is Niemeier's rank-$24$ class-number theorem itself. Path~(b)
is the explicit ADE rank-sum computation encoded in the displayed row
set. Path~(c) is the Schellekens-table count of the explicit row set.
This proves all claims. Throughout, the five objects are kept distinct:
$A$ is the VOA, $B(A)$ its bar coalgebra, $A^!$ its Koszul dual,
$A^i$ its dual coalgebra, and
$Z^{\mathrm{der}}_{\mathrm{ch}}(A)=\ChirHoch^\bullet(A,A)$ its derived
center; the theorem compares their invariants, not the objects
themselves.
\end{proof}

\begin{proposition}[Uniform weight-two threshold across the three
holomorphic \texorpdfstring{$c=24$}{c=24} machineries]
\label{prop:schellekens-weight-two-threshold}
\ClaimStatusProvedHere
On the ambient~$\mathfrak{S}_{24}^{\mathrm{hol}}$ of
Theorem~\ref{thm:schellekens-structured-subset}, the minimal uniform
threshold at which the Schellekens row, and therefore the structured
subset position of~$A$, is determined is the conformal-weight threshold
\[
w_\ast = 2.
\]
The same threshold is seen in each of the three machineries:
\begin{enumerate}[label=\textup{(\roman*)},leftmargin=*]
\item \textup{(}bar--cobar lane\textup{)}
 the weight-$\le2$ truncation of the ordered bar differential records
 the weight-one OPE residues and the first weight-two correction, hence
 reconstructs~$V_1$ and separates the Monster row~$V_1=0$ from the
 other~$70$ rows;
\item \textup{(}chiral Hochschild lane\textup{)}
 Theorem~H
 \textup{(}Theorem~\ref{thm:hochschild-concentration-E1} together with
 Corollary~\ref{cor:hochschild-averaging-symmetric}\textup{)}
 forces $\ChirHoch^\bullet(A)$ into degrees~$\{0,1,2\}$, and the
 weight-$\le2$ part detects exactly the same $V_1$/Monster dichotomy;
\item \textup{(}derived-center lane\textup{)}
 the evaluation pairing on
 $Z^{\mathrm{der}}_{\mathrm{ch}}(A)=\ChirHoch^\bullet(A,A)$ restricted
 to conformal weights~$\le2$ is read on the same underlying
 cochain complex as the Hochschild lane, so the threshold agrees
 tautologically with~\textup{(ii)}.
\end{enumerate}
In particular, the threshold statement is uniform across the three
machineries and is not an extra counting axiom.
\end{proposition}

\begin{proof}
Schellekens's classification is by~$V_1$ except for the unique Monster
row, whose first nonzero positive-weight datum lies in weight~$2$.
Therefore weight~$1$ is sufficient for the~$70$ non-Monster rows and
weight~$2$ is the minimal uniform threshold on the full family.

On the bar side, the reduced bar differential on weight-one generators
is the OPE residue Lie bracket, so the bar complex truncated at
conformal weight~$2$ already reconstructs~$V_1$ together with the
Monster exceptional case. On the Hochschild side,
Theorem~\ref{thm:hochschild-bar-cobar} identifies
$\ChirHoch^\bullet(A)$ with the bar-cobar coderivation complex, and
Theorem~\ref{thm:hochschild-concentration-E1} confines the live degrees
to~$\{0,1,2\}$; no higher cohomological threshold can intervene. On the
derived-center side,
$Z^{\mathrm{der}}_{\mathrm{ch}}(A)=\ChirHoch^\bullet(A,A)$ is the same
object with its $E_2$-structure retained, so the evaluation pairing on
weights~$\le2$ reads exactly the same data. Hence the common threshold
is~$w_\ast=2$ in all three lanes.
\end{proof}

\begin{remark}[Independent verification ledger for the numerical claims]
\label{rem:schellekens-numerical-ledger}
Each numerical claim used in
Theorem~\ref{thm:schellekens-structured-subset} and
Proposition~\ref{prop:schellekens-weight-two-threshold} has at least
three independent verification paths.
\begin{enumerate}[label=\textup{(\arabic*)},leftmargin=*]
\item \textup{(}$71 = 24 + 1 + 46$\textup{)}
 \begin{enumerate}[label=\textup{(\alph*)},leftmargin=*]
 \item Schellekens 1993, §3 and the numbered table on pp.~24--27 of the
  CERN/arXiv scan \textup{(}article pp.~177--182\textup{)}:
  rows~$0$ through~$70$ are
  explicit, and the row sets in
  Theorem~\ref{thm:schellekens-structured-subset}(i) partition them.
 \item Niemeier 1973, p.~142 \textup{(}abstract\textup{)}, and
  Conway--Sloane, Table~16.1, p.~407, give the~$24$ lattice rows;
  FLM88, Chapters~8--11, give the Monster singleton; subtraction from
  the total~$71$ leaves~$46$.
 \item M\"oller--Scheithauer 2023, Theorems~6.6, 6.7, and~6.9
  \textup{(}pp.~33--35, 40--41\textup{)} identify the nonzero
  $V_1$ rows, construct the nonzero Lie structures, and isolate
  the Monster sector, leaving the same~$46$-row non-lattice complement.
 \end{enumerate}
\item \textup{(}$24 = \operatorname{rank}(N) = \#\mathfrak{S}_{24}^{\mathrm{Niem}}$\textup{)}
 \begin{enumerate}[label=\textup{(\alph*)},leftmargin=*]
 \item Niemeier 1973, J. Number Theory~5, p.~142 \textup{(}abstract\textup{)},
  gives class number~$24$ for definite even quadratic forms of
  dimension~$24$ and discriminant~$1$.
 \item Conway--Sloane, Table~16.1, p.~407, together with the displayed
  ADE list above, yields the $23$ rootful simply-laced rank-$24$
  systems plus the rootless Leech case.
 \item The Schellekens-table intersection counts exactly the
  $24$ rows displayed in
  Theorem~\ref{thm:schellekens-structured-subset}(i).
 \end{enumerate}
\item \textup{(}central charge~$24$\textup{)}
 \begin{enumerate}[label=\textup{(\alph*)},leftmargin=*]
 \item It is the ambient hypothesis of Schellekens's classification.
 \item The lattice and Monster constructions compute
  $c=\operatorname{rank}(N)=24$ and $c(V^\natural)=24$ directly.
 \item The graded characters all begin with $q^{-1}$ and have modular
  weight~$0$, the holomorphic-$c=24$ normalization.
 \end{enumerate}
\item \textup{(}level~$1$ on the Niemeier rows\textup{)}
 \begin{enumerate}[label=\textup{(\alph*)},leftmargin=*]
 \item The lattice VOA construction produces simply-laced affine
  currents at level~$1$.
 \item Schellekens's numbered table on pp.~24--27 of the CERN/arXiv
  scan \textup{(}article pp.~177--182\textup{)} labels those rows
  with level-$1$ subscripts.
 \item The affine central-charge computation
  $c(\widehat{\fg}_{1})=\operatorname{rank}(\fg)$ on simply-laced
  factors recovers total central charge~$24$ exactly on those rows.
 \end{enumerate}
\item \textup{(}$w_\ast = 2$\textup{)}
 \begin{enumerate}[label=\textup{(\alph*)},leftmargin=*]
 \item Schellekens 1993, pp.~4, 23--24, classifies the nonzero cases by
  their weight-one Lie algebra and isolates the unique $N=0$ / Monster
  exceptional branch.
 \item The ordered bar differential on weight-one generators already
  reconstructs~$V_1$, and weight~$2$ is the first place the Monster
  branch differs.
 \item Theorem~H confines $\ChirHoch^\bullet$ to degrees~$\{0,1,2\}$,
  and the derived-center lane is the same cochain complex with its
  $E_2$-structure retained.
 \end{enumerate}
\end{enumerate}
\end{remark}

\begin{remark}[Cross-volume checks forced by the structured
decomposition]
\label{rem:schellekens-structured-cross-volume}
Theorem~\ref{thm:schellekens-structured-subset} has three immediate
cross-volume consequences.
\begin{enumerate}[label=\textup{(\arabic*)},leftmargin=*]
\item \textup{(}Vol~I Niemeier census\textup{)}
 the master-table row ``Niemeier $V_\Lambda$ (all~$24$)'' records only
 the coarse class-$G$ shadow slot. The theorem above supplies the
 missing structured subset behind that row.
\item \textup{(}Vol~III BKM / Fake Monster\textup{)}
 the number~$24$ appearing as Fake-Monster imaginary-root multiplicity
 is the Leech-rank input on the unique Leech/Niemeier branch; it is not
 the full holomorphic $c=24$ count by itself. The shared integer is a
 transport of the Leech slot, not an accidental slogan.
\item \textup{(}Vol~III categorical DT\textup{)}
 the portable statement on the CY$_3$ side is the low-weight threshold,
 not a new ``71'' count: the DT/Hochschild deformation data are read in
 degrees~$\le2$, mirroring Proposition~\ref{prop:schellekens-weight-two-threshold}.
\end{enumerate}
\end{remark}

\section{Coverage, propagation, and independent verification}
\label{sec:coverage-propagation}

The fingerprint classification is cross-volume: Vol I controls the
slots via the bar coalgebra (this chapter); Vol II reinterprets
the coset slot in the 3d HT gauge theory context; Vol III upgrades
the coset slot to $\mathrm{coset}_{FF}$-with-BKM-data on the CY
side. The critical-level structure (Theorem~\ref{thm:fifth-class-FF})
is the one slot in which all three volumes agree: the Feigin--Frenkel
centre appears uniformly in all three
\textup{(}\cite{FeiginFrenkel1992},~\cite{ArakawaFrenkel2017}\textup{)}.

\begin{remark}[Independent verification paths]
\label{rem:fingerprint-iv-paths}
The strengthening theorems and census surfaces in this chapter admit
disjoint verification sources:
\begin{enumerate}[label=\textup{(\alph*)},leftmargin=*]
\item Theorem~\ref{thm:schellekens-structured-subset} via three
 independent routes: Schellekens 1993, §3 and the numbered table on
 pp.~24--27 of the CERN/arXiv scan \textup{(}article pp.~177--182;
 holomorphic $c=24$ rows\textup{)}; Niemeier 1973, J. Number Theory~5, p.~142
 \textup{(}class number~$24$\textup{)} together with Conway--Sloane,
 Table~16.1, p.~407 \textup{(}rank-$24$ even unimodular lattices\textup{)};
 M\"oller--Scheithauer 2023, Theorems~6.6, 6.7, 6.9
 \textup{(}pp.~33--35, 40--41; generalised-deep-hole realization\textup{)}.
\item Theorem~\ref{thm:fifth-class-FF} via
 \cite{FeiginFrenkel1992} for the centre
 (algebraic) and \cite{BeilinsonDrinfeld1991} for the
 $\mathrm{Op}$-interpretation (geometric Langlands): two independent
 routes to the same infinite-dimensional structure.
\item Theorem~\ref{thm:pole-depth-independence} against the
Kac--Wakimoto classification of simple VOAs
\cite{KacWakimoto2004} and the Creutzig--Linshaw $W$-algebra
free-generation table \cite{CreutzigLinshaw2023}.
\item Theorem~\ref{thm:d-alg-r-max-bijection} against the
independently proved depth-gap trichotomy
\textup{(}Proposition~\ref{prop:depth-gap-trichotomy}\textup{)}.
\item Theorem~\ref{thm:DS-fingerprint-transport} against the
Kac--Roan--Wakimoto analysis of DS reductions
\cite{KacRoanWakimoto2003} and the Arakawa proof of DS quasi-isomorphism
\cite{Arakawa2007}.
\end{enumerate}
Compute-side harness at
\texttt{compute/tests/test\_infinite\_fingerprint.py}.
\end{remark}

\begin{remark}[Forward consequences and open direction]
\label{rem:fingerprint-forward-consequences}
The infinite fingerprint classification recasts several open
questions. (a) The Koszul dual of class~$\mathrm{FF}$ is again
class~$\mathrm{FF}$: the Feigin--Frenkel duality
$V_{-h^\vee}(\fg)^{\,!} \simeq \Fun\mathrm{Op}_{\fg^\vee}(D)$
is a genuine involution on the FF stratum. (b) Escalation of class
under DS reduction
(Theorem~\ref{thm:DS-fingerprint-transport}) is a one-way map
$\mathcal{F}_L\to\mathcal{F}_M$; the inverse problem (which class-$M$
fingerprints are DS images?) is open beyond the hook-type
Kac--Roan--Wakimoto range. (c) The fingerprint for elliptic and
toroidal chiral algebras requires an extension of slot~(5) to record
the KZB-connection data on the elliptic base and the Miki
$S_3$-symmetry on the toroidal base; this is open.
\end{remark}

% End of chapter.
